% =============================================================================
%  academic_paper_preamble.tex — Academic LaTeX paper preamble
%
%  Used by: academic_paper_latex_skill.md (working papers, journal
%           submissions, research documents).
%
%  HOW TO USE: copy the contents of this file into the top of the main .tex
%  file you produce. This file is a template, not a runtime include — do not
%  `\input{}` it. It begins with `\documentclass{article}` and must be the
%  literal start of the generated .tex file.
%
%  This is an ACADEMIC paper preamble. The document chrome (title,
%  section headers, links) follows academic conventions — black titles,
%  default-blue links via `blue!60!black`. No institutional branding: no
%  branded header, no logo footer, no colored title.
%
%  However, a color palette IS defined here for use in figures, TikZ
%  diagrams, and pgfplots charts. Figures often travel between an academic
%  working paper and a policy note as the same chart, so the same palette
%  is available in both contexts.
%
%  The color block below is this file's own definition — it is
%  self-contained and needs no external color reference. It ships with
%  placeholder colors: set the four base colors to your own before the
%  first real figure, and keep the color names stable so the rest of the
%  template keeps working.
%
%  Starter-package edition, mirrored from kit v6.11 (October 10, 2026). The
%  color block is the package's own. Record your own changes here.
% =============================================================================

\documentclass[11pt,letterpaper]{article}

% --- Page layout -------------------------------------------------------------
\usepackage[top=1in, bottom=1in, left=1in, right=1in]{geometry}

% --- Spacing -----------------------------------------------------------------
\usepackage{setspace}
\onehalfspacing
\raggedbottom

\usepackage{float}

% --- Encoding and fonts (pdflatex) -------------------------------------------
\usepackage[T1]{fontenc}
\usepackage[utf8]{inputenc}
\usepackage{textcomp}

% --- Math --------------------------------------------------------------------
\usepackage{amssymb,amsmath}

% --- Graphics and figures ----------------------------------------------------
\usepackage{graphicx}
\usepackage{subcaption}

\usepackage[most]{tcolorbox}

% --- Captions ----------------------------------------------------------------
\usepackage[
    justification=centering,
    labelfont=bf,
    labelsep=period,
    font=small,
    skip=6pt
]{caption}

% --- Tables ------------------------------------------------------------------
\usepackage{booktabs}
\usepackage{multirow}
\usepackage{tabularx}
\usepackage{longtable}
\usepackage{array}
\usepackage{calc}

% --- Colors — base xcolor for tables; the palette below for figures/TikZ -----
\usepackage[table,xcdraw]{xcolor}

% Color palette: define yours here.
% For use in figures, TikZ, and pgfplots. NOT for document chrome (titles,
% section headers, links) — those follow academic conventions.
% This template needs a small palette: three hues, each in a few shades, and
% grays. Set the four base colors below. The shades are derived from them, so
% the settings that use them keep working. The defaults are xcolor's generic
% colors: placeholders until you set your own. A fourth hue is added the same
% way.
\colorlet{BaseBlue}{blue!70!cyan!75!black}  % main hue: main series, points, lines
\colorlet{BaseGreen}{green!45!black}        % second hue: curves, contrasts
\colorlet{BaseYellow}{yellow!85!orange}     % warm hue: highlights
\colorlet{BaseGray}{black!45}               % labels, annotations

% Shades, derived from the bases. Keep these names stable.
\colorlet{PrimaryBlue}{BaseBlue}
\colorlet{DarkBlue}{BaseBlue!70!black}
\colorlet{SecondaryBlue}{BaseBlue!60}
\colorlet{SoftBlue}{BaseBlue!30}
\colorlet{LightBlue}{BaseBlue!18}
\colorlet{DarkGreen}{BaseGreen}
\colorlet{LightGreen}{BaseGreen!35}
\colorlet{WarmCream}{BaseYellow!20}
\colorlet{LabelGray}{BaseGray}
\colorlet{LightGray}{BaseGray!30}

% --- TikZ and pgfplots -------------------------------------------------------
\usepackage{tikz,pgfplots}
% pgfplotsset is optional — include only if the paper uses pgfplots charts
\pgfplotsset{%
   compat=newest,
   tick label style={font=\footnotesize},
   label style={font=\small},
   legend style={font=\small},
   axis x line = center,
   axis y line = center,
   every axis/.style={pin distance=1ex},
   trim axis left
}

% --- Footnotes ---------------------------------------------------------------
\usepackage[bottom]{footmisc}

% --- Multi-column ------------------------------------------------------------
\usepackage{multicol}

% --- Date formatting ---------------------------------------------------------
\usepackage[en-US]{datetime2}
\DTMnewdatestyle{monthyear}{%
  \renewcommand*{\DTMdisplaydate}[4]{\DTMenglishmonthname{##2} ##1}%
  \renewcommand*{\DTMDisplaydate}{\DTMdisplaydate}%
}

% --- Microtype ---------------------------------------------------------------
\usepackage{microtype}

% --- Line numbers (uncomment for draft review) -------------------------------
\usepackage{lineno}
%\linenumbers

% --- Bibliography ------------------------------------------------------------
\usepackage{natbib}
\bibliographystyle{aer}  % or chicago, plainnat, etc. — confirm with the user

% --- Hyperlinks (MUST be last package) ---------------------------------------
\usepackage{hyperref}
\hypersetup{
    colorlinks=true,
    linkcolor=blue!60!black,
    citecolor=blue!60!black,
    urlcolor=blue!60!black,
    pdftitle={Paper Title},
    pdfauthor={Author Names}
}

% =============================================================================
%  CUSTOM COMMANDS
% =============================================================================

\newcommand{\figurenote}[1]{%
    \vspace{4pt}%
    \begin{minipage}{0.95\linewidth}%
        \footnotesize \textit{Note:} #1%
    \end{minipage}%
}

\newcommand{\figuresource}[1]{%
    \vspace{2pt}%
    \begin{minipage}{0.95\linewidth}%
        \footnotesize \textit{Source:} #1%
    \end{minipage}%
}

\newcommand{\tablenote}[1]{%
    \vspace{4pt}%
    \begin{minipage}{0.95\linewidth}%
        \footnotesize \textit{Notes:} #1%
    \end{minipage}%
}

% Prevent overfull lines
\setlength{\emergencystretch}{3em}

